\documentclass{article}
\usepackage[backend=biber,natbib=true,style=alphabetic,maxbibnames=50]{biblatex}
\addbibresource{/home/nqbh/reference/bib.bib}
\usepackage[utf8]{vietnam}
\usepackage{tocloft}
\renewcommand{\cftsecleader}{\cftdotfill{\cftdotsep}}
\usepackage[colorlinks=true,linkcolor=blue,urlcolor=red,citecolor=magenta]{hyperref}
\usepackage{amsmath,amssymb,amsthm,enumitem,float,graphicx,mathtools,tikz}
\usetikzlibrary{angles,calc,intersections,matrix,patterns,quotes,shadings}
\allowdisplaybreaks
\newtheorem{assumption}{Assumption}
\newtheorem{baitoan}{Bài toán}
\newtheorem{cauhoi}{Câu hỏi}
\newtheorem{conjecture}{Conjecture}
\newtheorem{corollary}{Corollary}
\newtheorem{dangtoan}{Dạng toán}
\newtheorem{definition}{Definition}
\newtheorem{dinhly}{Định lý}
\newtheorem{dinhnghia}{Định nghĩa}
\newtheorem{example}{Example}
\newtheorem{ghichu}{Ghi chú}
\newtheorem{hequa}{Hệ quả}
\newtheorem{hypothesis}{Hypothesis}
\newtheorem{lemma}{Lemma}
\newtheorem{luuy}{Lưu ý}
\newtheorem{nhanxet}{Nhận xét}
\newtheorem{notation}{Notation}
\newtheorem{note}{Note}
\newtheorem{principle}{Principle}
\newtheorem{problem}{Problem}
\newtheorem{proposition}{Proposition}
\newtheorem{question}{Question}
\newtheorem{remark}{Remark}
\newtheorem{theorem}{Theorem}
\newtheorem{tip}{Tip}
\newtheorem{trick}{Trick}
\newtheorem{vidu}{Ví dụ}
\usepackage[left=1cm,right=1cm,top=5mm,bottom=5mm,footskip=4mm]{geometry}
\def\labelitemii{$\circ$}
\DeclareRobustCommand{\divby}{%
	\mathrel{\vbox{\baselineskip.65ex\lineskiplimit0pt\hbox{.}\hbox{.}\hbox{.}}}%
}
\setlist[itemize]{leftmargin=*}
\setlist[enumerate]{leftmargin=*}

\title{Vài Kiến Thức Cơ Bản, Tips, \& Tricks cho\\Olympic Toán Học Học Sinh {\it\&} Sinh Viên Toàn Quốc (VMC)}
\author{Nguyễn Quản Bá Hồng\footnote{A scientist- {\it\&} creative artist wannabe, a mathematics {\it\&} computer science lecturer of Department of Artificial Intelligence {\it\&} Data Science (AIDS), School of Technology (SOT), UMT Trường Đại học Quản lý {\it\&} Công nghệ TP.HCM, Hồ Chí Minh City, Việt Nam.\\E-mail: {\tt nguyenquanbahong@gmail.com, hong.nguyenquanba@umt.edu.vn}. Website: \url{https://nqbh.github.io/}. GitHub: \url{https://github.com/NQBH}.}}
\date{\today}

\begin{document}
\maketitle
\begin{abstract}
	Một số kiến thức, tip, \& trick để thi Olympic Toán Sinh Viên.
	
	Latest version:
	\begin{itemize}
		\item {\it Vài Kiến Thức Cơ Bản, Tips, \& Tricks cho Olympic Toán Học Học Sinh {\it\&} Sinh Viên Toàn Quốc (VMC)}.
		
		PDF: {\sc url}: \url{https://github.com/NQBH/advanced_STEM_beyond/blob/main/VMC/NQBH_VMC_lecture.pdf}.
		
		\TeX: {\sc url}: \url{https://github.com/NQBH/advanced_STEM_beyond/blob/main/VMC/NQBH_VMC_lecture.tex}.
	\end{itemize}
\end{abstract}
\tableofcontents

%------------------------------------------------------------------------------%

\section{Roadmap}

\paragraph{Styles.} Tài liệu này viết theo kiểu principle-based hơn là problem-based, i.e., tập trung vào các nguyên lý có tính liên kết khi phân tích \& giải 1 bài toán hơn là liệt kê các dạng toán khác nhau mà không liên kết chúng lại với nhau.

\paragraph{Generalization -- Tổng quát hóa.} Tác giả cho rằng việc quan trọng hơn việc học các tricks trong 1 lời giải cụ thể cho 1 bài toán cụ thể là việc hiểu lời giải hoặc chứng minh cụ thể đó hoạt động như thế nào, các lý luận chính (main arguments) là gì, \& có khi áp dụng được cho các bài toán hay dạng toán nào khác hay không. Khi đó, thông qua việc mở rộng phạm vi áp dụng (tựa tựa ``Bành Trướng Lãnh Địa'' trong Jujutsu Kaisen -- Chú thuật sư Hồi Chiến), các ý toán \& lập luận trong mỗi lời giải hoặc chứng minh cho 1 bài toán cụ thể sẽ có ý nghĩa hơn về mặt lâu dài.

%------------------------------------------------------------------------------%

\section{Algebra -- Đại Số}
Gồm 2 phần: Đại số tuyến tính (Linear Algebra) \& Đại số trừu tượng (Abstract Algebra). Olympic Toán Sinh Viên thường nghiêng về Đại số tuyến tính nhiều hơn.

\paragraph{Notation -- Ký hiệu.} Ta sử dụng ký hiệu $\mathbb{F}$ để chỉ 1 {\it trường} (field). Với sinh viên các ngành không phải Toán \& cận Toán thì chỉ cần quan tâm trường hợp $\mathbb{F} = \mathbb{R}$, còn các ngành cận Toán như Khoa học Máy tính, Công nghệ thông tin, Vật lý, etc., thì chỉ cần quan tâm các trường hợp $\mathbb{F}\in\{\mathbb{R},\mathbb{C}\}$, với sinh viên các ngành Toán, Toán--Tin, Toán--Cơ--Tin, nên quan tâm $\mathbb{F}\in\{\mathbb{R},\mathbb{C},\mathbb{Z}_p\mbox{ với } p \mbox{ nguyên tố}\}$ vì $\mathbb{Z}_p$ với $p$ nguyên tố, để làm các bài Đa thức \& Số học trong đề Đại số.

Với 1 ma trận $A\in M_{m\times n}(\mathbb{F})$, gọi $r_i(A),c_j(A)$ lần lượt là dòng (row) thứ $i\in[m]$, cột (column) thứ $j\in[n]$ của ma trận $A$.

%------------------------------------------------------------------------------%

\subsection{Some basic properties of matrices -- Vài tính chất cơ bản của ma trận}

\begin{remark}
	Ma trận nói chung không có tính chất giao hoán, i.e., $AB$ có thể khác $BA$, nhưng có tính chất kết hợp, i.e., $(AB)C = A(BC)$.
\end{remark}

\begin{remark}
	Khi 2 ma trận $A,B$ giao hoán với nhau, i.e., $AB = BA$ thì các hằng đẳng thức (đại số) đáng nhớ sẽ áp dụng được cho $A,B$, e.g.,
	\begin{align*}
		(A + B)^2 &= (A + B)(A + B) = A^2 + AB + BA + B^2 = A^2 + AB + AB + B^2 = A^2 + 2AB + B^2,\\
		(A - B)^2 &= (A - B)(A - B) = A^2 - AB - BA + B^2 = A^2 - AB - AB + B^2 = A^2 - 2AB + B^2. 
	\end{align*}
	Chú ý đẳng thức $AB = BA$ đã hàm ý luôn việc $A,B\in M_n(\mathbb{R})$ (i.e., 2 ma trận $A,B$ phải vuông \& cùng cấp). Thật vậy, giả sử $A\in  M_{m\times n}(\mathbb{F}),B\in M_{p\times q}(\mathbb{F})$, để tích $AB$ xác định thì $n = p$, \& để tích $BA$ xác định thì $q = m$. Khi đó $AB\in M_m(\mathbb{F}),BA\in M_n(\mathbb{F})$, đẳng thức $AB = BA$ suy ra $m = n$. Vậy $A,B\in M_n(\mathbb{R})$.
\end{remark}
Tương tự, ta có kết quả sau:

\begin{lemma}[Các hằng đẳng thức cho 2 ma trận giao hoán]
	Cho 2 ma trận $A,B\in M_n(\mathbb{C})$ giao hoán, i.e., $AB = BA$, \& 1 hằng đẳng thức đại số 2 biến có dạng $f(x,y) = 0$. Khi đó $f(A,B) = 0_n$.
	
	Nói riêng, ta có công thức nhị thức Newton sau cho 2 ma trận $A,B$:
	\begin{align*}
		(A + B)^n &= \sum_{i=0}^n C_n^iA^{n - i}B^i,\ \forall n\in\mathbb{N}^\star,\\
		(A - B)^n &= \sum_{i=0}^n (-1)^iC_n^iA^{n - i}B^i,\ \forall n\in\mathbb{N}^\star,\\
		A^n - B^n &= (A - B)\sum_{i=0}^{n-1} A^{n - 1 - i}B^i,\ \forall n\in\mathbb{N}^\star\\
		A^{2n + 1} - B^{2n + 1} &= (A + B)\sum_{i=0}^{2n} (-1)^iA^{2n - i}B^i,\ \forall n\in\mathbb{N}^\star
	\end{align*}
\end{lemma}

%------------------------------------------------------------------------------%

\subsection{Rank \& determinant of a matrix -- Hạng \& định thức của ma trận}

\begin{tip}
	Nếu bài Đại số cho 1 ma trận $4\times4$, thường là có tham số $a$ hoặc $m$, cố gắng tìm cấu trúc đặc biệt của ma trận đó, e.g., tổng các hàng{\tt/}cột, tổ hợp tuyến tính dễ thấy của các hàng{\tt/}cột.
\end{tip}
E.g.:

\begin{baitoan}[VMC2025A, Đại số, 1.]
	Cho $a\in\mathbb{R}$,
	\begin{equation*}
		M(a) = \begin{pmatrix}
			a & 2023 & 2024 & 2025\\2025 & 2024 & 2023 & a\\2024 & 2025 & a & 2023\\2023 & a & 2025 & 2024
		\end{pmatrix}
	\end{equation*}
	\item(a) Tính $\operatorname{rank}M(2022)$.
	\item(b) Tính $\det M(a)$ theo $a$.
	\item(c) Xét hệ phương trình tuyến tính thuần nhất $M(a)X = {\bf0}$, trong đó $X = (x_1,x_2,x_3,x_4)^\top$ là vector cột gồm 4 tọa độ $x_1,x_2,x_3,x_4\in\mathbb{R}$. Tìm số chiều của không gian nghiệm của hệ phương trình tuyến tính thuần nhất $M(2022)X = {\bf0}$. Khi nào hệ phương trình tuyến tính thuần nhất $M(a)X = {\bf0}$ chỉ có hữu hạn nghiệm?
\end{baitoan}

\begin{proof}[Phân tích]
	Đập ngay vào mắt, $M(a)$ có mỗi dòng \& mỗi cột đều chứa đúng 4 số $a,2023,2024,2025$ nên nếu thực hiện phép biến đổi sơ cấp 3 trên dòng{\tt/}cột bằng cách gán tổng của 4 dòng{\tt/}cột cho 1 dòng, e.g., dòng cuối $r_4$, thì sẽ thu được:
	\begin{equation*}
		r_4\coloneqq r_1 + r_2 + r_3 + r_4 = (a + 2023 + 2024 + 2025)(1,1,1,1),
	\end{equation*}
	nên nếu $a = -(2023 + 2024 + 2025)$ thì dòng $r_4$ này bằng $(0,0,0,0)$, nên đa thức đặc trưng $p_A(x)$ sẽ có 1 nghiệm là $x = -(2023 + 2024 + 2025) = -6072$.
\end{proof}

\begin{lemma}
	Nếu tồn tại 1 tổ hợp tuyến tính của các hàng hoặc các cột của 1 ma trận vuông $A\in M_n(\mathbb{C})$ thì $\det A = 0$.
\end{lemma}

\begin{lemma}
	Nếu $A\in M_n(\mathbb{R})$ giao hoán với mọi ma trận $B\in M_n(\mathbb{R})$ thì $A$ có dạng $A = cI_n$ với $c\in\mathbb{R}$.
\end{lemma}

\begin{lemma}
	Cho $A = {\rm diag}(a_1,a_2,\ldots,a_n),B = {\rm diag}(b_1,b_2,\ldots,b_n)$ thì
	\begin{align*}
		AB &= {\rm diag}(a_1b_1,a_2b_2,\ldots,a_nb_n),\\
		A^m &= {\rm diag}(a_1^m,a_2^m,\ldots,a_n^m),\ \forall m\in\mathbb{N}.
	\end{align*}
\end{lemma}

\begin{tip}
	Nếu bài Đại số yêu cầu tính $\det(A - \alpha I_n)$ thì phải xét đa thức đặc trưng $p_A(t) = \det(A - tI_n)$.
\end{tip}

%------------------------------------------------------------------------------%

\subsection{Trace of a matrix -- Vết ma trận}
Khái niệm {\it vết} (trace) chỉ được xác định với ma trận vuông, ma trận chữ nhật không vuông thì không có khái niệm vết.

\begin{dinhnghia}[Vết của ma trận]
	{\rm Vết} của 1 ma trận vuông $A = (a_{ij})_{i,j=1}^n\in M_n(\mathbb{F})$ được đặt bởi
	\begin{equation*}
		\operatorname{tr}(A) = \sum_{i=1}^n a_{ii}\in\mathbb{F}.
	\end{equation*}
\end{dinhnghia}
See, e.g., \href{https://en.wikipedia.org/wiki/Trace_(linear_algebra)}{Wikipedia{\tt/}trace (linear algebra)}.

\begin{definition}[Trace of a matrix]
	The {\rm trace} of a square matrix $A = (a_{ij})_{i,j=1}^n\in M_n(\mathbb{F})$, denoted by $\operatorname{tr}(A)$, is the sum of
\end{definition}

\paragraph{Lấy vết 2 vế của 1 phương trình ma trận.} Với 1 phương trình ma trận có dạng $f(A) = g(A)\in M_n(\mathbb{F})$

%------------------------------------------------------------------------------%

\subsection{Diagonalizability of matrices -- Tính chéo hóa được của ma trận}

\begin{dinhly}[Tính chéo hóa của ma trận]
	Nếu ma trận vuông $A\in M_n(\mathbb{R})$ có $n$ giá trị riêng phân biệt đôi một thì $A$ chéo hóa được.
\end{dinhly}

\begin{dinhly}[\cite{Hoa_linear_algebra}, Thm. 18.3]
	Cho $\{v_i\}_{i=1}^r$ là các vector riêng của $\varphi$ ứng với $r$ giá trị riêng $\{\lambda_i\}_{i=1}^r$ đôi một khác nhau. Khi đó $\{v_i\}_{i=1}^r$ độc lập tuyến tính.
\end{dinhly}

\begin{dinhnghia}
	1 toán tử tuyến tính được gọi là {\rm chéo hóa được} nếu nó có 1 ma trận biểu diễn là ma trận đường chéo. 1 ma trận vuông $A$ được gọi là {\rm chéo hóa được} nếu $A$ đồng dạng với ma trận đường chéo, i.e., tồn tại ma trận khả nghịch $P$ để $P^{-1}AP$ là 1 ma trận đường chéo.
\end{dinhnghia}

\begin{dinhly}[\cite{Hoa_linear_algebra}, Thm. 18.4]
	Cho $\varphi$ là 1 toán tử tuyến tính của không gian vector $V$ chiều $n$. Khi đó
	\item(i) $\varphi$ chéo hóa được khi \& chỉ khi $V$ có 1 cơ sở gồm các vector riêng.
	\item(ii) Nếu $\varphi$ có $n$ giá trị riêng khác nhau thì $\varphi$ chéo hóa được.
\end{dinhly}

\begin{baitoan}[VMC2025ĐSA2B4, \cite{VMS_VMC2025}, p. 21]
	Giả sử ban đầu học sinh các trường phổ thông ở Quận 1 chỉ mua bánh mì của Công ty A. Sau đó Công ty B cũng tham gia vào việc cung cấp bánh mì cho thị trường này. Nghiên cứu cho thấy sau mỗi quý (3 tháng) có $30\%$ khách hàng của Công ty A chuyển sang dùng sản phẩm của Công ty B, số còn lại vẫn dùng của Công ty A. Đồng thời có $40\%$ khách hàng của Công ty B chuyển sang dùng của Công ty A, số còn lại vẫn dùng của Công ty B.
	\item(a) Sau $18$ tháng kể từ khi Công ty B tham gia vào thị trường, tính tỷ lệ khách hàng (trên tổng số học sinh của toàn Quận 1) mua bánh mỳ của mỗi công ty tương ứng.
	\item(b) Giám đốc Công ty A quyết định sẽ ngừng kinh doanh mặt hàng bánh mì nếu tỷ lệ khách hàng của họ (trên tổng số) ít hơn $50\%$. Hỏi với số liệu như trên, Công ty A có thời điểm nào cần quyết định ngừng kinh doanh mặt hàng này không? Tại sao?
\end{baitoan}

\begin{proof}[Giải]
	Gọi $a_n,b_n$ lần lượt là tỷ lệ khách hàng trên tổng số ở chu kỳ thứ $n\in\mathbb{N}$, với $a_0 = 1,b_0 = 0$ (ban đầu chỉ có công ty A bán, công ty B chưa tham gia), lập được hệ phương trình truy hồi tuyến tính:
	\begin{equation*}
		\left\{\begin{split}
			a_{n+1} &= 0.7a_n + 0.4b_n,\\
			b_{n+1} &= 0.3a_n + 0.6b_n.
		\end{split}\right.
	\end{equation*}
	Đặt
	\begin{equation*}
		A = \begin{pmatrix}
			0.7 & 0.4\\0.3 & 0.6
		\end{pmatrix},
	\end{equation*}
	thì hệ phương trình truy hồi tuyến tính trên trở thành
	\begin{equation*}
		\begin{pmatrix}
			a_{n+1}\\b_{n+1}
		\end{pmatrix} = A\begin{pmatrix}
			a_n\\b_n
		\end{pmatrix}.
	\end{equation*}
	Chéo hóa ma trận $A$:
	\begin{equation*}
		A = \begin{pmatrix}
			\frac{4}{3} & -1\\1 & 1
		\end{pmatrix}\begin{pmatrix}
			1 & 0\\0 & \frac{3}{10}
		\end{pmatrix}\begin{pmatrix}
			\frac{4}{3} & -1\\1 & 1
		\end{pmatrix}^{-1}.
	\end{equation*}
	Lặp bằng hệ phương trình truy hồi tuyến tính (hoặc dễ dàng chứng minh bằng quy nạp):
	\begin{equation*}
		\begin{pmatrix}
			a_n\\b_n
		\end{pmatrix} = A\begin{pmatrix}
			a_{n-1}\\b_{n-1}
		\end{pmatrix} = A^2\begin{pmatrix}
			a_{n-2}\\b_{n-2}
		\end{pmatrix} = \cdots = A^n\begin{pmatrix}
			a_0\\b_0
		\end{pmatrix} = \begin{pmatrix}
			\frac{4}{3} & -1\\1 & 1
		\end{pmatrix}\begin{pmatrix}
			1 & 0\\0 & \left(\frac{3}{10}\right)^n
		\end{pmatrix}\begin{pmatrix}
			\frac{4}{3} & -1\\1 & 1
		\end{pmatrix}^{-1}\begin{pmatrix}
			1\\0
		\end{pmatrix} = \begin{pmatrix}
			\frac{4}{7} + \frac{3}{7}\left(\frac{3}{10}\right)^n\\\frac{3}{7} - \frac{3}{7}\left(\frac{3}{10}\right)^n
		\end{pmatrix},\ \forall n\in\mathbb{N}.
	\end{equation*}
	Sau 18 tháng $= 6$ quý, tỷ lệ khách hàng sử dụng sản phẩm 2 Công ty A, B tương ứng là:
	\begin{equation*}
		\begin{pmatrix}
			a_6\\b_6
		\end{pmatrix} = \begin{pmatrix}
			\frac{4}{7} + \frac{3}{7}\left(\frac{3}{10}\right)^6\\\frac{3}{7} - \frac{3}{7}\left(\frac{3}{10}\right)^6
		\end{pmatrix} = \begin{pmatrix}
			\frac{4002187}{7000000}\vspace{1mm}\\\frac{2997813}{7000000}.
		\end{pmatrix}
	\end{equation*}
	\item(b) Có dãy $a_n = \frac{4}{7} + \frac{3}{7}\left(\frac{3}{10}\right)^n$ giảm \& $\lim_{n\to\infty} a_n = \frac{4}{7} > 0.5$, nên tỷ lệ khách hàng sử dụng sản phẩm của Công ty A luôn $> \frac{4}{7} > \frac{3.5}{7} = 0.5$, nên Công ty A không cần phải thay đổi mặt hàng.
\end{proof}

\begin{remark}
	Bài toán trên là 1 trường hợp cụ thể của bài toán chuyển đổi trạng thái thông qua 1 hệ phương trình truy hồi tuyến tính được mô tả như sau: Giả sử có $n\in\mathbb{N}_{\ge2}$ đối tượng $\{A_i\}_{i=1}^n$ (agents, ở bài toán trên là 2 công ty A, B), mỗi đối tượng $A_i$ quan tâm đến 1 đại lượng $f_i(t_j)$ tại mỗi thời điểm $t_j\in[0,\infty)$, $\forall i\in[n]$, $\forall j\in\mathbb{N}$. Trạng thái ban đầu được cho trước: $(f_1(t_0),\ldots,f_n(t_0)) = (\alpha_1,\ldots,\alpha_n)$, \& quá trình chuyển trạng thái từ mốc thời gian $t_n$ sang $t_{n+1}$ được mô tả bởi hệ phương trình truy hồi tuyến tính sau:
	\begin{equation*}
		f_i(t_{n+1}) = \sum_{j=1}^n a_{ij}f_j(t_n),\ \forall i\in[n],\ \forall n\in\mathbb{N},
	\end{equation*}
	viết cụ thể:
	\begin{equation*}
		\left\{\begin{split}
			f_1(t_{n+1}) &= \sum_{j=1}^n a_{1j}f_j(t_n) = a_{11}f_1(t_n) + a_{12}f_2(t_n) + \cdots + a_{1n}f_n(t_n),\\
			f_2(t_{n+1}) &= \sum_{j=1}^n a_{2j}f_j(t_n) = a_{21}f_1(t_n) + a_{22}f_2(t_n) + \cdots + a_{2n}f_n(t_n),\\
			&\ldots\\
			f_n(t_{n+1}) &= \sum_{j=1}^n a_{nj}f_j(t_n) = a_{n1}f_1(t_n) + a_{n2}f_2(t_n) + \cdots + a_{nn}f_n(t_n),
		\end{split}\right.
	\end{equation*}
	hoặc dưới dạng ma trận nhân vector:
	\begin{equation*}
		\begin{pmatrix}
			f_1(t_{n+1})\\f_2(t_{n+1})\\\vdots\\f_n(t_{n+1})
		\end{pmatrix} = \begin{pmatrix}
			a_{11} & \cdots & a_{1n}\\a_{21} & \ddots & a_{2n}\\\vdots & \ddots & \vdots\\a_{n1} & \cdots & a_{nn}.
		\end{pmatrix}
	\end{equation*}
	Đặt $A = (a_{ij})_{i,j=1}^n$ \& ${\bf f}(t_{n+1}) = (f_1(t_{n+1}),f_2(t_{n+1}),\ldots,f_n(t_{n+1}))^\top$ thì
	\begin{equation*}
		{\bf f}(t_{n+1}) = A{\bf f}(t_n),\ \forall n\in\mathbb{N}.
	\end{equation*}
	Lặp:
	\begin{equation*}
		{\bf f}(t_m) = A{\bf f}(t_{m-1}) = A^2{\bf f}(t_{m-2}) = \cdots = A^m{\bf f}(t_0) = A^m\begin{pmatrix}
			\alpha_1\\\vdots\\\alpha_n
		\end{pmatrix},\ \forall m\in\mathbb{N}.
	\end{equation*}
	Đến đây ta cần tính $A^m$. Có 2 trường hợp:
	\begin{enumerate}
		\item {\bf Trường hợp $A$ chéo hóa được.} Tìm được ma trận $P\in M_n(\mathbb{R})$ khả nghịch thỏa $P^{-1}AP = \operatorname{diag}(\lambda_1(A),\ldots,\lambda_n(A))$ với $\{\lambda_i(A)\}_{i=1}^n$ là $n$ trị riêng (có thể gồm nghiệm riêng bội) của $A$, khi đó
		\begin{align*}
			A^m &= P\operatorname{diag}(\lambda_1^m(A),\ldots,\lambda_n^m(A))P^{-1},\ \forall m\in\mathbb{N},\\
			{\bf f}(t_m) &= P\operatorname{diag}(\lambda_1^m(A),\ldots,\lambda_n^m(A))P^{-1}\begin{pmatrix}
				\alpha_1\\\vdots\\\alpha_n
			\end{pmatrix},\ \forall m\in\mathbb{N}.
		\end{align*}
		\item {\bf Trường hợp $A$ không chéo hóa được.} Hy vọng $n$ đủ nhỏ hoặc có vài tính chất đơn giản để đơn giản được tính toán. Còn nếu không, trong trường hợp tổng quát, thì phải sử dụng binary exponentiation for square matrices bên Competitive Programming \& lập trình thì mới mong cứu nổi, see, e.g.
		
		NQBH. {\it Matrix Multiplication \& Fast Doubling Techniques in Competitive Programming}. \url{https://github.com/NQBH/advanced_STEM_beyond/blob/main/OLP_ICPC/matrix_multiplication/NQBH_matrix_multiplication.pdf}.
	\end{enumerate}
	
\end{remark}

%------------------------------------------------------------------------------%

\subsection{Matrix equation -- Phương trình ma trận}
Đề Olympic Toán Sinh viên thường xét các dạng toán sau, với $0_n = (0)_{i,j=1}^n$ là ma trận không kích thước $n\times n$:
\begin{enumerate}
	\item {\it Có tồn tại hay không 1 ma trận $A\in M_n(\mathbb{R})$ thỏa phương trình ma trận $f(A) = 0_n$?}
	
	Nếu tồn tại thì chỉ cần chỉ ra 1 ví dụ cho nghiệm ma trận $A$ thỏa mãn $f(A) = 0_n$ là xong. Còn nếu không tồn tại thì phải chứng minh $f(A)\ne0_n$, $\forall A\in M_n(\mathbb{R})$, \& thường sử dụng phương pháp phản chứng, dùng các tính chất quen thuộc của ma trận như các tính chất của det, trace, $\ldots$.
	\item {\it Chứng minh không tồn tại ma trận $A\in M_n(\mathbb{R})$ thỏa phương trình ma trận $f(A) = 0_n$.}
	
	I.e., vế sau của bài toán 1.
	\item {\it Tìm tất cả các ma trận $A\in M_n(\mathbb{R})$ thỏa mãn phương trình ma trận $f(A) = 0_n$.}
	
	Bài toán này yêu cầu giải phương trình ma trận để tìm tất cả các nghiệm ma trận.
\end{enumerate}

\begin{baitoan}[\cite{VMS_VMC2025}, UIT-HCM]
	\item(a) Có tồn tại hay không ma trận $A\in M_3(\mathbb{R})$ thỏa ${\rm tr}\,A = 0$, $A^2 + A^\top = I_3$?
	\item(b) Tìm tất cả các ma trận thực $X\in M_2(\mathbb{R})$ thỏa
	\begin{equation*}
		X^3 - 3X^2 + \begin{pmatrix}
			2 & 2\\2 & 2
		\end{pmatrix} = 0_2.
	\end{equation*}
\end{baitoan}

\begin{remark}[Phân tích ý (b)]
	\item(i) Phương trình ma trận của (b) có chứa 1 ma trận $2\times2$ cụ thể:
	\begin{equation*}
		A = \begin{pmatrix}
			2 & 2\\2 & 2
		\end{pmatrix}
	\end{equation*}
	Nhận thấy $\det A = 2\times2 - 2\times2 = 0$, nên ta có thể lấy $\det$ 2 vế để tận dụng tính chất $\det A = 0$. Hơn nữa,
	\begin{equation*}
		A - 4I_2 = \begin{pmatrix}
			2 & 2\\2 & 2
		\end{pmatrix} - 4\begin{pmatrix}
			1 & 0 \\0 & 1
		\end{pmatrix} = \begin{pmatrix}
			2 & -2\\-2 & 2
		\end{pmatrix}
	\end{equation*}
	cũng có $\det(A - 4I_2) = 2\times2 - (-2)\times(-2) = 0$ nên cũng có thể tận dụng tính chất này khi lấy $\det$ 2 vế, đương nhiên là sau khi thêm bớt số hạng $4I_n$ hợp lý.
	\item(ii) Nếu gặp các phương trình ma trận $f(X) = 0_n$ tương tự, với ẩn ma trận $X\in M_n(\mathbb{C})$, ta nên thử tìm các tính chất đặc biệt của các ma trận cụ thể trong phương trình ma trận đã cho, e.g., ma trận đó có định thức hoặc vết đặc biệt.
\end{remark}

\begin{tip}
	Để giải phương trình ma trận có dạng $f(A,A^\top) = 0_n$ với biến $A\in M_n(\mathbb{C})$, cố gắng biểu diễn $A^\top$ theo $A$, i.e., tìm hàm $g(x)$ để $A^\top = g(A)$, sau đó thay vào phương trình ma trận ban đầu để được $f(A,g(A)) = 0_n$, thường phương trình này sẽ là phương trình đa thức, cố gắng suy ra đa thức tối tiểu của $A$.
\end{tip}

\begin{baitoan}
	1 tổng quát hóa của ý (b) bài toán trên có dạng: Giải phương trình ma trận
	\begin{equation*}
		p(A,B)\coloneqq\sum_{i=1}^n a_iA^i + B = 0_n
	\end{equation*}
	với biến ma trận $A\in M_n(\mathbb{R})$, còn ma trận $B\in M_n(\mathbb{R})$ được cho trước (với vai trò như hệ số tự do trong 1 đa thức 1 biến $p\in\mathbb{R}[x]$).
\end{baitoan}

%------------------------------------------------------------------------------%

\subsection{Polynomial -- Đa thức}

\begin{baitoan}[VMC2025ĐSA5, \cite{VMS_VMC2025}, p. 22]
	\item(a) Xét đa thức $Q(x) = x^2 - 1$. Chứng minh rằng có vô số đa thức $P(x)$ bậc 2 với hệ số bậc cao nhất bằng $1$, sao cho đa thức $P(Q(x))$ có $4$ nghiệm thực phân biệt \& chúng lập thành 1 cấp số cộng.
	\item(b) Xét đa thức $T(x) = x^3 - 3x$. Chứng minh rằng không tồn tại đa thức $S(x)$ có bậc $3$ sao cho đa thức $S(T(x))$ có $9$ nghiệm thực phân biệt \& chúng lập thành 1 cấp số cộng.
\end{baitoan}

\begin{baitoan}
	Mở rộng bài toán VMC2025ĐSA5.
\end{baitoan}


%------------------------------------------------------------------------------%

\section{Analysis}

%------------------------------------------------------------------------------%

\subsection{Sequences -- Dãy số}
Cho 1 dãy số $\{a_n\}_{n=1}^\infty$. Đề Olympic Toán Sinh viên thường xét các dạng toán sau:
\begin{enumerate}
	\item Dãy số $\{a_n\}_{n=1}^\infty$ có hội tụ hay không?
	\item Chứng minh dãy số $\{a_n\}_{n=1}^\infty$ hội tụ hoặc phân kỳ.
	\item Tính giới hạn $\lim_{n\to\infty} a_n$.
\end{enumerate}

\begin{theorem}
	Nếu 1 dãy tăng ngặt (hoặc không giảm) \& bị chặn trên thì hội tụ. Nếu 1 dãy giảm ngặt (hoặc không tăng) thì hội tụ.
\end{theorem}

\begin{tip}
	Nếu dãy số có số hạng $u_n$ có công thức tổng quát là 1 tích, thì có thể dùng hàm $ln$ để biến tích thành tổng bằng cách đặt $v_n\coloneqq\ln u_n$.
\end{tip}

\begin{tip}[Kỹ thuật làm trội{\tt/}giảm]
	Nếu công thức có xuất hiện các số vô tỷ, hoặc không chính phương, có thể thay bằng số chính phương để tiện phân tích nhân tử. E.g., B1 bảng A Giải tích 2025, có thể làm trội như sau:
	\begin{equation*}
		a_n = \prod_{i=1}^n \left(1 - \frac{\pi}{(i + 1)^2}\right) > \prod_{i=1}^n \left(1 - \frac{4}{(i + 1)^2}\right) = \prod_{i=1}^n \left(1 - \frac{1}{i + 1}\right)\left(1 + \frac{1}{i + 1}\right)
	\end{equation*}
	vì $\pi$ là số vô tỷ, khá khó chịu, nên làm trội bởi $\pi < 4$ với $4$ là 1 số chính phương gần nhất.
\end{tip}

\begin{tip}
	Nếu cần xét dấu đạo hàm $f'(x)$, có thể phải sử dụng bất đẳng thức Cauchy--Schwartz:
	\begin{equation*}
		\left(\sum_{i=1}^n a_i^2\right)\left(\sum_{i=1}^n b_i^2\right)\ge\left(\sum_{i=1}^n a_ib_i\right)^2,\ \forall a_i,b_i\in\mathbb{R},\ \forall i\in[n].
	\end{equation*}
	Bất đẳng thức Cauchy--Schwartz là hệ quả trực tiếp của hằng đẳng thức Lagrange:
	\begin{equation*}
		\left(\sum_{i=1}^n a_i^2\right)\left(\sum_{i=1}^n b_i^2\right) = \left(\sum_{i=1}^n a_ib_i\right)^2 + \sum_{1\le i < j\le n} (a_ib_j - a_jb_i)^2\ge\left(\sum_{i=1}^n a_ib_i\right)^2.
	\end{equation*}
	Nhưng cần phải nhớ kỹ công thức hằng đẳng thức Lagrange, không thì cứ sử dụng bất đẳng thức Cauchy--Schwartz.
\end{tip}

\begin{example}
	
\end{example}

\begin{tip}
	Nếu dính bài cần đánh giá $\cos$ thì sử dụng bất đẳng thức
	\begin{equation*}
		\cos x\ge1 - \frac{x^2}{2},\ \forall x\ge0.
	\end{equation*}
\end{tip}

\begin{dinhly}[Darboux]
	Nếu đạo hàm $f'$ của 1 hàm số $f$ nhận 2 giá trị $a,b\in\mathbb{R}$ thì $f'$ cũng nhận tất cả các giá trị thuộc $[a,b]$.
\end{dinhly}
I.e.,
\begin{equation*}
	a,b\in{\rm Im}(f')\Leftrightarrow[a,b]\subset{\rm Im}(f').
\end{equation*}

\begin{tip}[Kỹ thuật liên tục hóa bài toán rời rạc]
	Đối với các bài yêu cầu tính giới hạn của dãy số $\{a_n\}_{n=1}^\infty$, nên chuyển từ phiên bản ``rời rạc'' thành phiên bản ``liên tục'' bằng cách tìm hàm số $f\in C(\mathbb{R})$ để $a_{n+1} = f(a_n)$, $\forall n\in\mathbb{N}^\star$ hoặc từ 1 số hạng nào đó trở đi, ngoại trừ vài số hạng đầu, để khảo sát tính chất của hàm số $f$, từ đó suy ra tính chất của dãy số $\{a_n\}_{n=1}^\infty$. Lợi thế của kỹ thuật liên tục hóa là ở phiên bản liên tục, có các công cụ giải tích mạnh như đạo hàm, tích phân, khảo sát hàm số, etc.
\end{tip}

%------------------------------------------------------------------------------%

\subsection{Continuous function -- Hàm số liên tục}
{\it Intuition.} Hàm số liên tục là hàm mà khi ta vẽ đồ thị của hàm đó, thì ta không được phép nhấc cây viết lên. Nếu nhấc lên thì hàm số sẽ bị gián đoạn (discontinuities) tại điểm nhấc lên đó, nên hàm số đó không liên tục (discontinuous function).

Gọi $F_{\rm elm}(\mathbb{F})$ là tập hợp các {\it hàm số sơ cấp (elementary function)} trên trường $\mathbb{F}$, see, e.g., \href{https://en.wikipedia.org/wiki/Elementary_function}{Wikipedia{\tt/}elementary functions}.



\begin{dinhnghia}[Hàm số liên tục tại điểm]
	
\end{dinhnghia}

\begin{dinhnghia}[Hàm số liên tục trên miền]
	
\end{dinhnghia}

\begin{lemma}[Định lý kẹp cho giới hạn hàm số]
	Cho $f,g:D\to\mathbb{R}$ là 2 hàm số thỏa mãn $|f(x)|\le g(x)$, $\forall x\in D$, $x_0\in D$, \& $\lim_{x\to x_0} g(x) = 0$. Khi đó $\lim_{x\to x_0} f(x) = 0$.
\end{lemma}

\begin{tip}
	Để chứng minh hàm số $f(x)$ không liên tục tại $x = x_0$, có thể tìm 2 dãy số $\{x_n\}_{n=1}^\infty,\{y_n\}_{n=1}^\infty$ cùng hội tụ về $x_0$ nhưng giới hạn của 2 dãy $\{f(x_n)\}_{n=1}^\infty,\{f(y_n)\}_{n=1}^\infty$lại khác nhau, i.e.,
	\begin{equation*}
		\left\{\begin{split}
			\lim_{n\to\infty} x_n &= \lim_{n\to\infty} y_n = x_0,\\
			\lim_{n\to\infty} f(x_n)&\ne\lim_{n\to\infty} f(y_n).
		\end{split}\right.		
	\end{equation*}
	Khi đó, theo định nghĩa liên tục điểm của hàm số, $f$ không liên tục (gián đoạn) tại $x = x_0$.
\end{tip}

%------------------------------------------------------------------------------%

\subsection{Derivatives of differentiable functions -- Các đạo hàm của các hàm số khả vi}

\begin{tip}
	Đối với các hàm số không phải là hàm sơ cấp, đặc biệt là hàm ghép lại của các hàm sơ cấp, i.e.,
	\begin{equation*}
		f(x) = \sum_{i=1}^{n-1} f_i(x){\bf1}_{(a_i,a_{i+1}]}(x) = \left\{\begin{split}
			&f_1(x)&&\mbox{if } x\in(a_1,a_2],\\
			&f_2(x)&&\mbox{if } x\in(a_2,a_3],\\
			&\ldots\\
			&f_n(x)&&\mbox{if } x\in(a_{n-1},a_n],
		\end{split}\right.
	\end{equation*}
	với $f_i\in F_{\rm elm}(\mathbb{R})$, $\forall i\in[n]$, $a_1,a_n\in\overline{\mathbb{R}}$, $a_i\in\mathbb{R},\ \forall i\in[[2,n-1]]$ thỏa $a_i < a_{i+1}$, $\forall i\in[n - 1]$, thì nên dùng định nghĩa của đạo hàm thông qua giới hạn để tính $f'(x)$.
\end{tip}


%------------------------------------------------------------------------------%

\subsection{Integral of integrable functions -- Tích phân của các hàm khả tích}

%------------------------------------------------------------------------------%

\section{Combinatorics -- Tổ Hợp}
Các bài thi Tổ hợp thì hên xui, nên bắt đầu thử các trường hợp nhỏ của $n$ nếu đề yêu cầu giải quyết cho trường hợp $n$ tổng quát, để tạo ``cảm giác'' đối với cấu trúc của bài toán Tổ hợp.

Nếu bài toán Tổ hợp yêu cầu tìm các giá trị lớn nhất hoặc nhỏ nhất, thì thường đó sẽ là bài toán Tổ hợp Cực biên, nên thử các cấu hình ở biên để có thêm thông tin.

%------------------------------------------------------------------------------%

\section{Number Theory -- Số Học}

%------------------------------------------------------------------------------%

\section{Miscellaneous -- Linh Tinh}
Trong thời đại AI hiện nay, đa số học sinh, sinh viên sẽ sử dụng AI rất nhiều. Tác giả hy vọng các bài viết truyền thống kiểu này sẽ vẫn giữ vững giá trị.

%------------------------------------------------------------------------------%

\paragraph{Acknowledgments.} Tác giả NQBH xin cảm ơn thầy{\tt/}anh {\sc Lê Phúc Lữ} đã tặng quyển sách \cite{Vinh_Lu_Olympic_Toan} để làm 1 nguồn tài liệu tham khảo quý báo trong việc biên soạn các tài liệu Olympic, cảm ơn sinh viên {\sc Phan Vĩnh Tiến}, cùng CLB Học thuật STAC UMT đã biên soạn các đề thi chọn đội tuyển cấp trường được chia sẽ trên Group Hướng tới Olympic Toán, cảm ơn trường UMT đã tạo điều kiện tốt nhất để lớp bồi dưỡng Toán, dù với thời lượng ngắn ngủi, đã đạt được chất lượng như mong đợi.

%------------------------------------------------------------------------------%

\printbibliography[heading=bibintoc]

\end{document}